\nwfilename{nw/xunit.nw}\nwbegindocs{0}% -*- mode: poly-noweb; noweb-code-mode: sml-mode; -*-% ===> this file was generated automatically by noweave --- better not edit it
%%
%% xunit.nw
%% 
%% Made by Alex Nelson <pqnelson@gmail.com>
%% Login   <alex@lisp>
%% 
%% Started on  2026-05-04T16:56:51-0700
%% Last update 2026-05-04T16:56:51-0700
%% 

\section{Unit testing framework}

\begin{node}
How do we know our code works as expected? One solution is to
\emph{prove} the code works according to some specification and some
formal semantics of the programming language. This is hard,
time-consuming, and brittle.

Another solution is to test code against a handful of
examples. This is easier, more robust, and acts as a form of
documentation (want to know how to use the code? Look at some
examples, in the form of tests). Further, testing gives us immediate
feedback about whether our expected use case works or
not. (Admittedly, in the static mode of presentation, this dynamic
feedback loop is missing.)
\end{node}

\begin{node}[Overview of design]
We define each \define{Test Case} as a function which checks one
scenario for a ``system under test'' (usually a specific
function). The typical test cases are either (1) simple ``check the
actual output matches the expected output'' tests, or (2) check
specific exceptions are raised in illegal situations.

A generic test cases has four ``phases'' or ``steps'':
\begin{enumerate}
\item \textsc{Setup} the test fixture [all the data needed to describe
  the ``before executing the system under test''];
\item \textsc{Exercise} the system under test (i.e., run the code you
  want to test against the test fixture);
\item \textsc{Verify} the ``actual output'' matches [in some
  appropriate sense] the ``expected output''
\item \textsc{Teardown} the test fixtures.
\end{enumerate}
\SML/ written in a functional style handles the teardown for us
automatically, sparing us one step of montony.

The xUnit framework uses a \define{Assertion Method} to verify the
actual output matches the expected output. 

A collection of tests form a \define{Test Suite}, which is itself a
test. In particular, a test suite may contain a nested test
suite. Every test is either a test case or a test suite.

The \define{Test Runner} takes tests, executes them, constructing
\define{Test Results} containing the outcome of each test along with
other metadata (like the amount of time it took to perform). We keep
things as simple as possible, and manually register tests with the
test runner.

We then have a \define{Test Reporter} take test results produced by
the runner, then prints a summary to the screen (or otherwise produces
some kind of artifact, like an XML file). Some people advocate saving
an XML file for each time you run unit tests, other people do not. We
are agnostic, but want to recognize this is a source of change (so we
abstract it away into its own module).
\end{node}

\begin{node}[Assertions]
We have a simple module to handle assertions for equality, or
asserting that a certain condition is satisfied. User-specific
messages are supplied for failure reporting.

When the assertion is not satisfied, a custom {\Tt{}\nwlinkedidentq{Assert.Fail}{NWsR2q7-4gQMdZ-1}\nwendquote}
exception is propagated. Later, the test runner will handle these
exceptions differently than any other: assertion failures are recorded
as ``failed tests'' whereas all other exceptional situations are
recorded as ``unexpected exception raised'' as its result.
\end{node}

\nwenddocs{}\nwbegincode{1}\sublabel{NWsR2q7-4gQMdZ-1}\nwmargintag{{\nwtagstyle{}\subpageref{NWsR2q7-4gQMdZ-1}}}\moddef{Assert.sml~{\nwtagstyle{}\subpageref{NWsR2q7-4gQMdZ-1}}}\endmoddef\nwstartdeflinemarkup\nwenddeflinemarkup
structure Assert = struct
exception Fail of string;

(* that : bool -> string -> unit *)
fun that cond msg =
  if cond then ()
  else raise Fail msg;

(* eq : ''a -> ''a -> string -> unit *)
fun eq expected actual msg =
  that (expected = actual) msg;
end;
\nwindexdefn{\nwixident{Assert.Fail}}{Assert.Fail}{NWsR2q7-4gQMdZ-1}\nwindexdefn{\nwixident{Assert.that}}{Assert.that}{NWsR2q7-4gQMdZ-1}\nwindexdefn{\nwixident{Assert.eq}}{Assert.eq}{NWsR2q7-4gQMdZ-1}\eatline
\nwnotused{Assert.sml}\nwidentdefs{\\{{\nwixident{Assert.eq}}{Assert.eq}}\\{{\nwixident{Assert.Fail}}{Assert.Fail}}\\{{\nwixident{Assert.that}}{Assert.that}}}\nwendcode{}\nwbegindocs{2}\nwdocspar
\begin{node}[Test]
We begin with the data structure for a test case and a test suite.

A test case has a name and a ``test method'' which the runner
executes. We construct a test case with {\Tt{}\nwlinkedidentq{Test.mk}{NWsR2q7-3fh2dK-1}\nwendquote} (since {\Tt{}mk\nwendquote} is
the standard constructor prefix/``name component'').

A test suite has a name and a collection of tests. We construct a test
suite with the {\Tt{}\nwlinkedidentq{Test.suite}{NWsR2q7-3fh2dK-1}\nwendquote} method.

Executing a test suite or test case is handled with the {\Tt{}\nwlinkedidentq{Test.run}{NWsR2q7-1bpV8D-1}\nwendquote}
method. It will produce a {\Tt{}\nwlinkedidentq{Test.Result.t}{NWsR2q7-oQiQW-1}\nwendquote} object storing the outcome
and metadata.
\end{node}

\nwenddocs{}\nwbegincode{3}\sublabel{NWsR2q7-1Fzsdc-1}\nwmargintag{{\nwtagstyle{}\subpageref{NWsR2q7-1Fzsdc-1}}}\moddef{Test.sig~{\nwtagstyle{}\subpageref{NWsR2q7-1Fzsdc-1}}}\endmoddef\nwstartdeflinemarkup\nwenddeflinemarkup
signature Test = sig
  type t;
  structure Result : sig
              \LA{}Specification for test result~{\nwtagstyle{}\subpageref{NWsR2q7-1ETFk0-1}}\RA{}
            end;
  val suite : string -> t list -> t;
  val mk : string -> (unit -> unit) -> t;
  val run : t -> Result.t;
end;
\nwnotused{Test.sig}\nwendcode{}\nwbegindocs{4}\nwdocspar

\nwenddocs{}\nwbegincode{5}\sublabel{NWsR2q7-3ucTCC-1}\nwmargintag{{\nwtagstyle{}\subpageref{NWsR2q7-3ucTCC-1}}}\moddef{Test.sml~{\nwtagstyle{}\subpageref{NWsR2q7-3ucTCC-1}}}\endmoddef\nwstartdeflinemarkup\nwenddeflinemarkup
structure Test :> Test = struct

structure Result = struct
\LA{}Implementation of test result~{\nwtagstyle{}\subpageref{NWsR2q7-oQiQW-1}}\RA{}
end;

\LA{}Constructors for tests~{\nwtagstyle{}\subpageref{NWsR2q7-3fh2dK-1}}\RA{}

\LA{}Run a test case or suite~{\nwtagstyle{}\subpageref{NWsR2q7-1bpV8D-1}}\RA{}
end;
\nwnotused{Test.sml}\nwendcode{}\nwbegindocs{6}\nwdocspar

\begin{node}
The constructors for a test case and test suite are rather straightforward.
\end{node}

\nwenddocs{}\nwbegincode{7}\sublabel{NWsR2q7-3fh2dK-1}\nwmargintag{{\nwtagstyle{}\subpageref{NWsR2q7-3fh2dK-1}}}\moddef{Constructors for tests~{\nwtagstyle{}\subpageref{NWsR2q7-3fh2dK-1}}}\endmoddef\nwstartdeflinemarkup\nwusesondefline{\\{NWsR2q7-3ucTCC-1}}\nwenddeflinemarkup
datatype t = Case of string * (unit -> unit)
           | Suite of string * (t list);

fun mk name method = Case(name, method);

fun suite name tests = Suite(name, tests);
\nwindexdefn{\nwixident{Test.t}}{Test.t}{NWsR2q7-3fh2dK-1}\nwindexdefn{\nwixident{Test.mk}}{Test.mk}{NWsR2q7-3fh2dK-1}\nwindexdefn{\nwixident{Test.suite}}{Test.suite}{NWsR2q7-3fh2dK-1}\eatline
\nwused{\\{NWsR2q7-3ucTCC-1}}\nwidentdefs{\\{{\nwixident{Test.mk}}{Test.mk}}\\{{\nwixident{Test.suite}}{Test.suite}}\\{{\nwixident{Test.t}}{Test.t}}}\nwendcode{}\nwbegindocs{8}\nwdocspar
\begin{node}
We abstract away the outcome of the test result, since more may be
added later. Some people want to allow ``skipping'' tests and wish to
add ``skipped'' as an outcome. The outcomes we will report are:
Success (the test's outcome was as expected), ``failure'' (the test's
outcome was not as expected, so we need the failure message), and
``error'' (an unexpected exception was raised, which we store in the
outcome).

The test result for a test case includes the test name, the amount of
time elapsed while running the test as measured against a ``clock'',
and the test outcome.

The results for a test suite are similar, but instead of a single test
outcome we have a collection of test results (for all the tests
included in the suite).
\end{node}

\nwenddocs{}\nwbegincode{9}\sublabel{NWsR2q7-1ETFk0-1}\nwmargintag{{\nwtagstyle{}\subpageref{NWsR2q7-1ETFk0-1}}}\moddef{Specification for test result~{\nwtagstyle{}\subpageref{NWsR2q7-1ETFk0-1}}}\endmoddef\nwstartdeflinemarkup\nwusesondefline{\\{NWsR2q7-1Fzsdc-1}}\nwprevnextdefs{\relax}{NWsR2q7-1ETFk0-2}\nwenddeflinemarkup
type t;
\nwalsodefined{\\{NWsR2q7-1ETFk0-2}\\{NWsR2q7-1ETFk0-3}\\{NWsR2q7-1ETFk0-4}\\{NWsR2q7-1ETFk0-5}\\{NWsR2q7-1ETFk0-6}\\{NWsR2q7-1ETFk0-7}\\{NWsR2q7-1ETFk0-8}\\{NWsR2q7-1ETFk0-9}}\nwused{\\{NWsR2q7-1Fzsdc-1}}\nwendcode{}\nwbegindocs{10}\nwdocspar

\nwenddocs{}\nwbegincode{11}\sublabel{NWsR2q7-oQiQW-1}\nwmargintag{{\nwtagstyle{}\subpageref{NWsR2q7-oQiQW-1}}}\moddef{Implementation of test result~{\nwtagstyle{}\subpageref{NWsR2q7-oQiQW-1}}}\endmoddef\nwstartdeflinemarkup\nwusesondefline{\\{NWsR2q7-3ucTCC-1}}\nwprevnextdefs{\relax}{NWsR2q7-oQiQW-2}\nwenddeflinemarkup
datatype Outcome = Success
                 | Failure of string
                 | Error of exn;

datatype t = Case of string * Time.time * Outcome
           | Suite of string * Time.time * (t list);
\nwindexdefn{\nwixident{Test.Result.t}}{Test.Result.t}{NWsR2q7-oQiQW-1}\eatline
\nwalsodefined{\\{NWsR2q7-oQiQW-2}\\{NWsR2q7-oQiQW-3}\\{NWsR2q7-oQiQW-4}\\{NWsR2q7-oQiQW-5}\\{NWsR2q7-oQiQW-6}\\{NWsR2q7-oQiQW-7}\\{NWsR2q7-oQiQW-8}\\{NWsR2q7-oQiQW-9}}\nwused{\\{NWsR2q7-3ucTCC-1}}\nwidentdefs{\\{{\nwixident{Test.Result.t}}{Test.Result.t}}}\nwendcode{}\nwbegindocs{12}\nwdocspar
\begin{node}
How do we actually ``construct'' a test result from a test? Well, we
run it. For a test case, this amounts to ``clocking'' the test and
recording the outcome. We'll delegate this work to {\Tt{}\nwlinkedidentq{Test.Result.for{\_}case}{NWsR2q7-oQiQW-2}\nwendquote}.

For a test suite, we have a bunch of tests to run. We'll just lazily
run all these tests.
\end{node}

\nwenddocs{}\nwbegincode{13}\sublabel{NWsR2q7-1bpV8D-1}\nwmargintag{{\nwtagstyle{}\subpageref{NWsR2q7-1bpV8D-1}}}\moddef{Run a test case or suite~{\nwtagstyle{}\subpageref{NWsR2q7-1bpV8D-1}}}\endmoddef\nwstartdeflinemarkup\nwusesondefline{\\{NWsR2q7-3ucTCC-1}}\nwenddeflinemarkup
(* run : \nwlinkedidentc{Test.t}{NWsR2q7-3fh2dK-1} -> \nwlinkedidentc{Test.Result.t}{NWsR2q7-oQiQW-1} *)
fun run (Case (name, method)) = Result.for_case name method
  | run (Suite (name, tests)) =
      Result.for_suite name (fn () => map run tests);
\nwindexdefn{\nwixident{Test.run}}{Test.run}{NWsR2q7-1bpV8D-1}\eatline
\nwused{\\{NWsR2q7-3ucTCC-1}}\nwidentdefs{\\{{\nwixident{Test.run}}{Test.run}}}\nwidentuses{\\{{\nwixident{Test.Result.t}}{Test.Result.t}}\\{{\nwixident{Test.t}}{Test.t}}}\nwindexuse{\nwixident{Test.Result.t}}{Test.Result.t}{NWsR2q7-1bpV8D-1}\nwindexuse{\nwixident{Test.t}}{Test.t}{NWsR2q7-1bpV8D-1}\nwendcode{}\nwbegindocs{14}\nwdocspar
\begin{node}
Given the name (as a string) and the four-phase test procedure as a
function, we:
\begin{enumerate}
\item Record the start time
\item Execute the test
\item In the default case: record the time interval since we started
  the test, along with the name, and that the test outcome was a {\Tt{}Success\nwendquote}.
\item If an {\Tt{}\nwlinkedidentq{Assert.Fail}{NWsR2q7-4gQMdZ-1}\ msg\nwendquote} was raised, then we record the time
  interval since we started, along with the name, and that the outcome
  was a {\Tt{}Failure\nwendquote} with the given {\Tt{}msg\nwendquote}.
\item If any other exception {\Tt{}e\nwendquote} was raised, then we record the time
  interval since we started, along with the test case's name, and that
  the outcome was an {\Tt{}Error\ e\nwendquote}.
\end{enumerate}
\end{node}

\nwenddocs{}\nwbegincode{15}\sublabel{NWsR2q7-1ETFk0-2}\nwmargintag{{\nwtagstyle{}\subpageref{NWsR2q7-1ETFk0-2}}}\moddef{Specification for test result~{\nwtagstyle{}\subpageref{NWsR2q7-1ETFk0-1}}}\plusendmoddef\nwstartdeflinemarkup\nwusesondefline{\\{NWsR2q7-1Fzsdc-1}}\nwprevnextdefs{NWsR2q7-1ETFk0-1}{NWsR2q7-1ETFk0-3}\nwenddeflinemarkup
val for_case : string -> (unit -> unit) -> t;
\nwused{\\{NWsR2q7-1Fzsdc-1}}\nwendcode{}\nwbegindocs{16}\nwdocspar

\nwenddocs{}\nwbegincode{17}\sublabel{NWsR2q7-oQiQW-2}\nwmargintag{{\nwtagstyle{}\subpageref{NWsR2q7-oQiQW-2}}}\moddef{Implementation of test result~{\nwtagstyle{}\subpageref{NWsR2q7-oQiQW-1}}}\plusendmoddef\nwstartdeflinemarkup\nwusesondefline{\\{NWsR2q7-3ucTCC-1}}\nwprevnextdefs{NWsR2q7-oQiQW-1}{NWsR2q7-oQiQW-3}\nwenddeflinemarkup
fun for_case name assertion =
  let
    val start = Time.now();
  in
    (assertion();
     Case(name,
          (Time.-)(Time.now(), start),
          Success))
    handle \nwlinkedidentc{Assert.Fail}{NWsR2q7-4gQMdZ-1} msg =>
           Case (name,
                 (Time.-)(Time.now(), start),
                 Failure msg)
         | e => Case(name,
                     (Time.-)(Time.now(), start),
                     Error e)
  end;
\nwindexdefn{\nwixident{Test.Result.for{\_}case}}{Test.Result.for:uncase}{NWsR2q7-oQiQW-2}\eatline
\nwused{\\{NWsR2q7-3ucTCC-1}}\nwidentdefs{\\{{\nwixident{Test.Result.for{\_}case}}{Test.Result.for:uncase}}}\nwidentuses{\\{{\nwixident{Assert.Fail}}{Assert.Fail}}}\nwindexuse{\nwixident{Assert.Fail}}{Assert.Fail}{NWsR2q7-oQiQW-2}\nwendcode{}\nwbegindocs{18}\nwdocspar
\begin{node}
The test result for a test suite is determined by:
\begin{enumerate}
\item Clock the start time; then
\item Calculate the result of each test by recursively invoking
  {\Tt{}\nwlinkedidentq{Test.run}{NWsR2q7-1bpV8D-1}\nwendquote} on each test in the suite; then
\item Clock the end time, recording the duration of running the tests
  in the suite; then
\item Finally construct the test result object for the test suite.
\end{enumerate}
\end{node}

\nwenddocs{}\nwbegincode{19}\sublabel{NWsR2q7-1ETFk0-3}\nwmargintag{{\nwtagstyle{}\subpageref{NWsR2q7-1ETFk0-3}}}\moddef{Specification for test result~{\nwtagstyle{}\subpageref{NWsR2q7-1ETFk0-1}}}\plusendmoddef\nwstartdeflinemarkup\nwusesondefline{\\{NWsR2q7-1Fzsdc-1}}\nwprevnextdefs{NWsR2q7-1ETFk0-2}{NWsR2q7-1ETFk0-4}\nwenddeflinemarkup
val for_suite : string -> (unit -> (t list)) -> t;
\nwused{\\{NWsR2q7-1Fzsdc-1}}\nwendcode{}\nwbegindocs{20}\nwdocspar

\nwenddocs{}\nwbegincode{21}\sublabel{NWsR2q7-oQiQW-3}\nwmargintag{{\nwtagstyle{}\subpageref{NWsR2q7-oQiQW-3}}}\moddef{Implementation of test result~{\nwtagstyle{}\subpageref{NWsR2q7-oQiQW-1}}}\plusendmoddef\nwstartdeflinemarkup\nwusesondefline{\\{NWsR2q7-3ucTCC-1}}\nwprevnextdefs{NWsR2q7-oQiQW-2}{NWsR2q7-oQiQW-4}\nwenddeflinemarkup
fun for_suite name mk_results =
  let
    val start = Time.now();
    val results = mk_results ();
    val dt = (Time.-)(Time.now(), start);
  in
    Suite(name, dt, results)
  end;
\nwindexdefn{\nwixident{Test.Result.for{\_}suite}}{Test.Result.for:unsuite}{NWsR2q7-oQiQW-3}\eatline
\nwused{\\{NWsR2q7-3ucTCC-1}}\nwidentdefs{\\{{\nwixident{Test.Result.for{\_}suite}}{Test.Result.for:unsuite}}}\nwendcode{}\nwbegindocs{22}\nwdocspar
\begin{node}[Accessor methods for test result objects]
We have abstracted the test result internals away, so we won't have to
worry about them. We have standard ways to get the name of the test,
the failure message (if any---otherwise return the empty string), and
the exception message (if any---otherwise return the empty string).
\end{node}

\nwenddocs{}\nwbegincode{23}\sublabel{NWsR2q7-1ETFk0-4}\nwmargintag{{\nwtagstyle{}\subpageref{NWsR2q7-1ETFk0-4}}}\moddef{Specification for test result~{\nwtagstyle{}\subpageref{NWsR2q7-1ETFk0-1}}}\plusendmoddef\nwstartdeflinemarkup\nwusesondefline{\\{NWsR2q7-1Fzsdc-1}}\nwprevnextdefs{NWsR2q7-1ETFk0-3}{NWsR2q7-1ETFk0-5}\nwenddeflinemarkup
val name : t -> string;
val msg : t -> string;
val exn_message : t -> string;
\nwused{\\{NWsR2q7-1Fzsdc-1}}\nwendcode{}\nwbegindocs{24}\nwdocspar

\nwenddocs{}\nwbegincode{25}\sublabel{NWsR2q7-oQiQW-4}\nwmargintag{{\nwtagstyle{}\subpageref{NWsR2q7-oQiQW-4}}}\moddef{Implementation of test result~{\nwtagstyle{}\subpageref{NWsR2q7-oQiQW-1}}}\plusendmoddef\nwstartdeflinemarkup\nwusesondefline{\\{NWsR2q7-3ucTCC-1}}\nwprevnextdefs{NWsR2q7-oQiQW-3}{NWsR2q7-oQiQW-5}\nwenddeflinemarkup
(* name : \nwlinkedidentc{Test.Result.t}{NWsR2q7-oQiQW-1} -> string *)
fun name (Case (n,_,_)) = n
  | name (Suite (n,_,_)) = n;

(* msg : \nwlinkedidentc{Test.Result.t}{NWsR2q7-oQiQW-1} -> string *)
fun msg (Case (_,_,Failure msg)) = msg
  | msg _ = "";

(* exn_msg : \nwlinkedidentc{Test.Result.t}{NWsR2q7-oQiQW-1} -> string *)
fun exn_message (Case (_,_,Error e)) = exnMessage e
  | exn_message _ = ""; 
\nwindexdefn{\nwixident{Test.Result.name}}{Test.Result.name}{NWsR2q7-oQiQW-4}\nwindexdefn{\nwixident{Test.Result.msg}}{Test.Result.msg}{NWsR2q7-oQiQW-4}\nwindexdefn{\nwixident{Test.Result.exn{\_}message}}{Test.Result.exn:unmessage}{NWsR2q7-oQiQW-4}\eatline
\nwused{\\{NWsR2q7-3ucTCC-1}}\nwidentdefs{\\{{\nwixident{Test.Result.exn{\_}message}}{Test.Result.exn:unmessage}}\\{{\nwixident{Test.Result.msg}}{Test.Result.msg}}\\{{\nwixident{Test.Result.name}}{Test.Result.name}}}\nwidentuses{\\{{\nwixident{Test.Result.t}}{Test.Result.t}}}\nwindexuse{\nwixident{Test.Result.t}}{Test.Result.t}{NWsR2q7-oQiQW-4}\nwendcode{}\nwbegindocs{26}\nwdocspar
\begin{node}
For test cases, we can access the time it took to execute.

For test suites, there is some overhead by recursively invoking
{\Tt{}\nwlinkedidentq{Test.run}{NWsR2q7-1bpV8D-1}\nwendquote} lazily. If the user wants to omit this, they can
calculate the sum of time intervals it took to run each test case.
\end{node}

\nwenddocs{}\nwbegincode{27}\sublabel{NWsR2q7-1ETFk0-5}\nwmargintag{{\nwtagstyle{}\subpageref{NWsR2q7-1ETFk0-5}}}\moddef{Specification for test result~{\nwtagstyle{}\subpageref{NWsR2q7-1ETFk0-1}}}\plusendmoddef\nwstartdeflinemarkup\nwusesondefline{\\{NWsR2q7-1Fzsdc-1}}\nwprevnextdefs{NWsR2q7-1ETFk0-4}{NWsR2q7-1ETFk0-6}\nwenddeflinemarkup
(* realtime excludes the time it took a TestSuite to
   allocate infrastructure in memory. *)
val realtime : t -> Time.time;
val runtime : t -> Time.time;
\nwused{\\{NWsR2q7-1Fzsdc-1}}\nwendcode{}\nwbegindocs{28}\nwdocspar

\nwenddocs{}\nwbegincode{29}\sublabel{NWsR2q7-oQiQW-5}\nwmargintag{{\nwtagstyle{}\subpageref{NWsR2q7-oQiQW-5}}}\moddef{Implementation of test result~{\nwtagstyle{}\subpageref{NWsR2q7-oQiQW-1}}}\plusendmoddef\nwstartdeflinemarkup\nwusesondefline{\\{NWsR2q7-3ucTCC-1}}\nwprevnextdefs{NWsR2q7-oQiQW-4}{NWsR2q7-oQiQW-6}\nwenddeflinemarkup
(* runtime : t -> Time.time

How long did it take to run the test(s) and construct the
result(s)? *)
  fun runtime (Case (_,dt,_)) = dt
    | runtime (Suite (_,dt,_)) = dt;

  (* realtime : t -> Time.time

How long did it take just to run the test(s)? *)
  fun realtime (Case (_,dt,_)) = dt
    | realtime (Suite (_,_,[])) = Time.zeroTime
    | realtime (Suite (_,_,[x])) = realtime x
    | realtime (Suite (_,_,r::rs)) =
      foldl (fn (result,dt) =>
                (Time.+)(dt, realtime result))
            (realtime r)
            rs;
\nwindexdefn{\nwixident{Test.Result.runtime}}{Test.Result.runtime}{NWsR2q7-oQiQW-5}\nwindexdefn{\nwixident{Test.Result.realtime}}{Test.Result.realtime}{NWsR2q7-oQiQW-5}\eatline
\nwused{\\{NWsR2q7-3ucTCC-1}}\nwidentdefs{\\{{\nwixident{Test.Result.realtime}}{Test.Result.realtime}}\\{{\nwixident{Test.Result.runtime}}{Test.Result.runtime}}}\nwendcode{}\nwbegindocs{30}\nwdocspar
\begin{node}
When reporting test results for a test suite, we will need to access
the results for each test in the suite.
\end{node}

\nwenddocs{}\nwbegincode{31}\sublabel{NWsR2q7-1ETFk0-6}\nwmargintag{{\nwtagstyle{}\subpageref{NWsR2q7-1ETFk0-6}}}\moddef{Specification for test result~{\nwtagstyle{}\subpageref{NWsR2q7-1ETFk0-1}}}\plusendmoddef\nwstartdeflinemarkup\nwusesondefline{\\{NWsR2q7-1Fzsdc-1}}\nwprevnextdefs{NWsR2q7-1ETFk0-5}{NWsR2q7-1ETFk0-7}\nwenddeflinemarkup
val results : t -> t list;
\nwused{\\{NWsR2q7-1Fzsdc-1}}\nwendcode{}\nwbegindocs{32}\nwdocspar

\nwenddocs{}\nwbegincode{33}\sublabel{NWsR2q7-oQiQW-6}\nwmargintag{{\nwtagstyle{}\subpageref{NWsR2q7-oQiQW-6}}}\moddef{Implementation of test result~{\nwtagstyle{}\subpageref{NWsR2q7-oQiQW-1}}}\plusendmoddef\nwstartdeflinemarkup\nwusesondefline{\\{NWsR2q7-3ucTCC-1}}\nwprevnextdefs{NWsR2q7-oQiQW-5}{NWsR2q7-oQiQW-7}\nwenddeflinemarkup
fun results (Suite (_,_,rs)) = rs;
\nwindexdefn{\nwixident{Test.Result.results}}{Test.Result.results}{NWsR2q7-oQiQW-6}\eatline
\nwused{\\{NWsR2q7-3ucTCC-1}}\nwidentdefs{\\{{\nwixident{Test.Result.results}}{Test.Result.results}}}\nwendcode{}\nwbegindocs{34}\nwdocspar
\begin{node}
We can tabulate the number of successes, failures, errors, and total
number of tests. The logic for these functions are all rather similar.
\begin{itemize}
\item For a test case, if the outcome is a {\Tt{}Success\nwendquote} then return
  $1$, otherwise return $0$.
\item For a test suite, just count all the successes in each test in
  the suite.
\end{itemize}
Counting the total number of test cases just sums the total number of
successes, failures, and errors.
\end{node}

\nwenddocs{}\nwbegincode{35}\sublabel{NWsR2q7-1ETFk0-7}\nwmargintag{{\nwtagstyle{}\subpageref{NWsR2q7-1ETFk0-7}}}\moddef{Specification for test result~{\nwtagstyle{}\subpageref{NWsR2q7-1ETFk0-1}}}\plusendmoddef\nwstartdeflinemarkup\nwusesondefline{\\{NWsR2q7-1Fzsdc-1}}\nwprevnextdefs{NWsR2q7-1ETFk0-6}{NWsR2q7-1ETFk0-8}\nwenddeflinemarkup
val count_successes : t -> int;
val count_failures : t -> int;
val count_errors : t -> int;
val count_total : t -> int;
\nwused{\\{NWsR2q7-1Fzsdc-1}}\nwendcode{}\nwbegindocs{36}\nwdocspar

\nwenddocs{}\nwbegincode{37}\sublabel{NWsR2q7-oQiQW-7}\nwmargintag{{\nwtagstyle{}\subpageref{NWsR2q7-oQiQW-7}}}\moddef{Implementation of test result~{\nwtagstyle{}\subpageref{NWsR2q7-oQiQW-1}}}\plusendmoddef\nwstartdeflinemarkup\nwusesondefline{\\{NWsR2q7-3ucTCC-1}}\nwprevnextdefs{NWsR2q7-oQiQW-6}{NWsR2q7-oQiQW-8}\nwenddeflinemarkup
fun count_successes (Case (_,_,Success)) = 1
  | count_successes (Case _) = 0 
  | count_successes (Suite (_,_,outcomes)) =
      foldl (op +) 0 (map count_successes outcomes);

fun count_failures (Case (_,_,Failure _)) = 1
  | count_failures (Case _) = 0
  | count_failures (Suite (_,_,outcomes)) =
      foldl (op +) 0 (map count_failures outcomes);

fun count_errors (Case (_,_,Error _)) = 1
  | count_errors (Case _) = 0
  | count_errors (Suite (_,_,outcomes)) =
      foldl (op +) 0 (map count_errors outcomes);

fun count_total x =
  count_successes x + count_failures x + count_errors x;
\nwindexdefn{\nwixident{Test.Result.count{\_}total}}{Test.Result.count:untotal}{NWsR2q7-oQiQW-7}\nwindexdefn{\nwixident{Test.Result.count{\_}successes}}{Test.Result.count:unsuccesses}{NWsR2q7-oQiQW-7}\nwindexdefn{\nwixident{Test.Result.count{\_}failures}}{Test.Result.count:unfailures}{NWsR2q7-oQiQW-7}\nwindexdefn{\nwixident{Test.Result.count{\_}errors}}{Test.Result.count:unerrors}{NWsR2q7-oQiQW-7}\eatline
\nwused{\\{NWsR2q7-3ucTCC-1}}\nwidentdefs{\\{{\nwixident{Test.Result.count{\_}errors}}{Test.Result.count:unerrors}}\\{{\nwixident{Test.Result.count{\_}failures}}{Test.Result.count:unfailures}}\\{{\nwixident{Test.Result.count{\_}successes}}{Test.Result.count:unsuccesses}}\\{{\nwixident{Test.Result.count{\_}total}}{Test.Result.count:untotal}}}\nwendcode{}\nwbegindocs{38}\nwdocspar
\begin{node}[Predicate for test outcome]
We can say a test suite or a test case is a success if the {\Tt{}count{\_}successes\nwendquote} equals the {\Tt{}count{\_}total\nwendquote}.

Failure for a test case is rather obvious (the outcome was a
{\Tt{}Failure\nwendquote}). But what qualifies as a ``failure'' for a test suite? We
will treat \emph{any} reported failure as indicating the test suite
should qualify as a ``failure''.

Similarly, any error result in a test contained in a test suite
qualifies that test suite as a ``error''.
\end{node}

\nwenddocs{}\nwbegincode{39}\sublabel{NWsR2q7-1ETFk0-8}\nwmargintag{{\nwtagstyle{}\subpageref{NWsR2q7-1ETFk0-8}}}\moddef{Specification for test result~{\nwtagstyle{}\subpageref{NWsR2q7-1ETFk0-1}}}\plusendmoddef\nwstartdeflinemarkup\nwusesondefline{\\{NWsR2q7-1Fzsdc-1}}\nwprevnextdefs{NWsR2q7-1ETFk0-7}{NWsR2q7-1ETFk0-9}\nwenddeflinemarkup
val is_success : t -> bool;
val is_failure : t -> bool;
val is_error : t -> bool;
\nwused{\\{NWsR2q7-1Fzsdc-1}}\nwendcode{}\nwbegindocs{40}\nwdocspar

\nwenddocs{}\nwbegincode{41}\sublabel{NWsR2q7-oQiQW-8}\nwmargintag{{\nwtagstyle{}\subpageref{NWsR2q7-oQiQW-8}}}\moddef{Implementation of test result~{\nwtagstyle{}\subpageref{NWsR2q7-oQiQW-1}}}\plusendmoddef\nwstartdeflinemarkup\nwusesondefline{\\{NWsR2q7-3ucTCC-1}}\nwprevnextdefs{NWsR2q7-oQiQW-7}{NWsR2q7-oQiQW-9}\nwenddeflinemarkup
fun is_success x = (count_total x) = (count_successes x);

fun is_failure (Case (_,_, Failure _)) = true
  | is_failure (Case _) = false
  | is_failure (s as Suite _) = (count_failures s) > 0;

fun is_error (Case (_,_, Error _)) = true
  | is_error (Case _) = false
  | is_error (s as Suite _) = (count_errors s) > 0;
\nwindexdefn{\nwixident{Test.Result.is{\_}success}}{Test.Result.is:unsuccess}{NWsR2q7-oQiQW-8}\nwindexdefn{\nwixident{Test.Result.is{\_}failure}}{Test.Result.is:unfailure}{NWsR2q7-oQiQW-8}\nwindexdefn{\nwixident{Test.Result.is{\_}error}}{Test.Result.is:unerror}{NWsR2q7-oQiQW-8}\eatline
\nwused{\\{NWsR2q7-3ucTCC-1}}\nwidentdefs{\\{{\nwixident{Test.Result.is{\_}error}}{Test.Result.is:unerror}}\\{{\nwixident{Test.Result.is{\_}failure}}{Test.Result.is:unfailure}}\\{{\nwixident{Test.Result.is{\_}success}}{Test.Result.is:unsuccess}}}\nwendcode{}\nwbegindocs{42}\nwdocspar
\begin{node}[Predicates for test results]
Since the test result is an opaque data structure, we need to provide
predicates to indicate whether a test result is for a suite or a test
case. This will be necessary for reporting results.
\end{node}

\nwenddocs{}\nwbegincode{43}\sublabel{NWsR2q7-1ETFk0-9}\nwmargintag{{\nwtagstyle{}\subpageref{NWsR2q7-1ETFk0-9}}}\moddef{Specification for test result~{\nwtagstyle{}\subpageref{NWsR2q7-1ETFk0-1}}}\plusendmoddef\nwstartdeflinemarkup\nwusesondefline{\\{NWsR2q7-1Fzsdc-1}}\nwprevnextdefs{NWsR2q7-1ETFk0-8}{\relax}\nwenddeflinemarkup
val is_case : t -> bool;
val is_suite : t -> bool;
\nwused{\\{NWsR2q7-1Fzsdc-1}}\nwendcode{}\nwbegindocs{44}\nwdocspar

\nwenddocs{}\nwbegincode{45}\sublabel{NWsR2q7-oQiQW-9}\nwmargintag{{\nwtagstyle{}\subpageref{NWsR2q7-oQiQW-9}}}\moddef{Implementation of test result~{\nwtagstyle{}\subpageref{NWsR2q7-oQiQW-1}}}\plusendmoddef\nwstartdeflinemarkup\nwusesondefline{\\{NWsR2q7-3ucTCC-1}}\nwprevnextdefs{NWsR2q7-oQiQW-8}{\relax}\nwenddeflinemarkup
fun is_case (Case _) = true
  | is_case _ = false;

fun is_suite (Suite _) = true
  | is_suite _ = false;
\nwindexdefn{\nwixident{Test.Result.is{\_}case}}{Test.Result.is:uncase}{NWsR2q7-oQiQW-9}\nwindexdefn{\nwixident{Test.Result.is{\_}suite}}{Test.Result.is:unsuite}{NWsR2q7-oQiQW-9}\eatline
\nwused{\\{NWsR2q7-3ucTCC-1}}\nwidentdefs{\\{{\nwixident{Test.Result.is{\_}case}}{Test.Result.is:uncase}}\\{{\nwixident{Test.Result.is{\_}suite}}{Test.Result.is:unsuite}}}\nwendcode{}\nwbegindocs{46}\nwdocspar
\begin{node}[Test reporter specification]
So far, we have implemented the assertion methods, the test case and
suite structures, and the test result data structure. We just need a
test runner and test reporter modules. Let us now turn to the test reporter.

We will abstract away ``reporting'' as any \SML/ function producing a
unit type (which is what happens when we write to a file or print to
the screen). Ostensibly, the user may want to send an email (or post a
TikTok video, or page someone, or\dots), so we just abstract away
whatever ``reporting'' mechanism this could be.
\end{node}

\nwenddocs{}\nwbegincode{47}\sublabel{NWsR2q7-jZez2-1}\nwmargintag{{\nwtagstyle{}\subpageref{NWsR2q7-jZez2-1}}}\moddef{Reporter.sig~{\nwtagstyle{}\subpageref{NWsR2q7-jZez2-1}}}\endmoddef\nwstartdeflinemarkup\nwenddeflinemarkup
signature Reporter = sig
  val report : \nwlinkedidentc{Test.Result.t}{NWsR2q7-oQiQW-1} -> unit;
  val report_all : \nwlinkedidentc{Test.Result.t}{NWsR2q7-oQiQW-1} list -> unit;
end;
\nwnotused{Reporter.sig}\nwidentuses{\\{{\nwixident{Test.Result.t}}{Test.Result.t}}}\nwindexuse{\nwixident{Test.Result.t}}{Test.Result.t}{NWsR2q7-jZez2-1}\nwendcode{}\nwbegindocs{48}\nwdocspar

\begin{node}
I'm used the JUnit's report to the terminal, so I'm going to stick to
that. We will report the time interval in microseconds.
\end{node}

\nwenddocs{}\nwbegincode{49}\sublabel{NWsR2q7-bqErk-1}\nwmargintag{{\nwtagstyle{}\subpageref{NWsR2q7-bqErk-1}}}\moddef{JUnitTt.sml~{\nwtagstyle{}\subpageref{NWsR2q7-bqErk-1}}}\endmoddef\nwstartdeflinemarkup\nwenddeflinemarkup
(* 
Print to the terminal (hence the "-Tt" suffix) a summary of test
results in the style of JUnit.

If any failures or errors occur, print those out with a bit more
detail.
*)
structure JUnitTt :> Reporter = struct
structure Result = Test.Result;

fun interval_to_string (dt : Time.time) =
  (LargeInt.toString (Time.toMicroseconds dt))^"ms";

\LA{}Iteratively report JUnit-style report for a single outcome~{\nwtagstyle{}\subpageref{NWsR2q7-4wVoS-1}}\RA{}

val report = print o (report_iter "");

val report_all = app report;
end;
\nwnotused{JUnitTt.sml}\nwendcode{}\nwbegindocs{50}\nwdocspar

\begin{node}
How we report results depends on whether we are reporting for a test
case or a test suite:
\begin{itemize}
\item Reporting a test case:
\begin{enumerate}
\item If the test failed, we should print the path to the test, announce it failed, and give the failure message.
\item If the test err'd, we should do everything similarly to the failure case, except announce it was an ERROR and produce the exception message.
\item If the test succeeded, then don't do anything: we do not wish to pollute the terminal.
\end{enumerate}
\item For a test suite:
\begin{enumerate}
\item Write the test suite's full path to the screen.
\item Write the results for each test in the suite to the screen.
\item Write to the screen a summary line counting the total number of tests run, the number of failures, the number of errors, and the name of the test suite.
\end{enumerate}
\end{itemize}
\end{node}

\nwenddocs{}\nwbegincode{51}\sublabel{NWsR2q7-4wVoS-1}\nwmargintag{{\nwtagstyle{}\subpageref{NWsR2q7-4wVoS-1}}}\moddef{Iteratively report JUnit-style report for a single outcome~{\nwtagstyle{}\subpageref{NWsR2q7-4wVoS-1}}}\endmoddef\nwstartdeflinemarkup\nwusesondefline{\\{NWsR2q7-bqErk-1}}\nwenddeflinemarkup
fun report_iter p =
  let
    fun path "" s = s
      | path p s = p^"."^s; 
    fun for_case c =
      if Result.is_failure c
      then concat [path p (Result.name c),
                   " FAIL: ",
                   Result.msg c,
                   "\\n"]
      else if Result.is_error c
      then concat [path p (Result.name c),
                   " ERROR: ",
                   (Result.exn_message c),
                   "\\n"]
      else "";
    fun for_suite result =
      let
        val p2 = path p (Result.name result);
        val rs = \nwlinkedidentc{Test.Result.results}{NWsR2q7-oQiQW-6} result;
      in
        concat ["Running ", p2, "\\n",
                concat (map (report_iter p2) rs),
                "Tests run: ",
                Int.toString(Result.count_total result),
                ", Failures: ",
                Int.toString(Result.count_failures result),
                ", Errors: ",
                Int.toString(Result.count_errors result),
                ", Time elapsed: ",
                interval_to_string (Result.realtime result),
                " - in ",
                (Result.name result), "\\n"]
      end;
  in
    fn r => if \nwlinkedidentc{Test.Result.is_case}{NWsR2q7-oQiQW-9} r then for_case r
            else for_suite r
  end;
\nwused{\\{NWsR2q7-bqErk-1}}\nwidentuses{\\{{\nwixident{Test.Result.is{\_}case}}{Test.Result.is:uncase}}\\{{\nwixident{Test.Result.results}}{Test.Result.results}}}\nwindexuse{\nwixident{Test.Result.is{\_}case}}{Test.Result.is:uncase}{NWsR2q7-4wVoS-1}\nwindexuse{\nwixident{Test.Result.results}}{Test.Result.results}{NWsR2q7-4wVoS-1}\nwendcode{}\nwbegindocs{52}\nwdocspar

\begin{node}[Test runner]
Now we have arrived at the final stage of the unit testing framework:
we want to take a list of tests, execute them, and report the results
using some test reporter. We will simply create a ``parametrized module''
(\SML/ functor) which allows for any arbitrary test reporter to be used.

The test runner has a single method, which expects a list of tests to
run and report.
\end{node}

\nwenddocs{}\nwbegincode{53}\sublabel{NWsR2q7-4Peuwc-1}\nwmargintag{{\nwtagstyle{}\subpageref{NWsR2q7-4Peuwc-1}}}\moddef{TestRunner.sig~{\nwtagstyle{}\subpageref{NWsR2q7-4Peuwc-1}}}\endmoddef\nwstartdeflinemarkup\nwenddeflinemarkup
signature TestRunner = sig
  val run : \nwlinkedidentc{Test.t}{NWsR2q7-3fh2dK-1} list -> unit;
end;
\nwnotused{TestRunner.sig}\nwidentuses{\\{{\nwixident{Test.t}}{Test.t}}}\nwindexuse{\nwixident{Test.t}}{Test.t}{NWsR2q7-4Peuwc-1}\nwendcode{}\nwbegindocs{54}\nwdocspar

\begin{node}
The test runner:
\begin{enumerate}
\item Runs the tests to obtain the test results,
\item Reports the results,
\end{enumerate}
and that's it.
\end{node}

\nwenddocs{}\nwbegincode{55}\sublabel{NWsR2q7-2dKvEw-1}\nwmargintag{{\nwtagstyle{}\subpageref{NWsR2q7-2dKvEw-1}}}\moddef{tests/MkRunner.sml~{\nwtagstyle{}\subpageref{NWsR2q7-2dKvEw-1}}}\endmoddef\nwstartdeflinemarkup\nwenddeflinemarkup
(*
Using any test reporter, we can run the tests, then
report the results.
*)
functor MkRunner(Reporter : Reporter) :> TestRunner = struct
  fun run tests =
    let
      val results = map \nwlinkedidentc{Test.run}{NWsR2q7-1bpV8D-1} tests;
    in
      Reporter.report_all results;
      ()
    end;
end;
\nwnotused{tests/MkRunner.sml}\nwidentuses{\\{{\nwixident{Test.run}}{Test.run}}}\nwindexuse{\nwixident{Test.run}}{Test.run}{NWsR2q7-2dKvEw-1}\nwendcode{}\nwbegindocs{56}\nwdocspar

\begin{node}
The ``main'' function is rather straightforward for unit tests.

Working with MLton and Poly/ML, we also define the {\Tt{}main\nwendquote} function here
for future usage.

We also initialize the unit tests to be run with an empty test
suite. This way, when we register test suites, we will do it in a
uniform manner (to aid readability).
\end{node}

\nwenddocs{}\nwbegincode{57}\sublabel{NWsR2q7-3MeQDa-1}\nwmargintag{{\nwtagstyle{}\subpageref{NWsR2q7-3MeQDa-1}}}\moddef{tests/main.sml~{\nwtagstyle{}\subpageref{NWsR2q7-3MeQDa-1}}}\endmoddef\nwstartdeflinemarkup\nwenddeflinemarkup
(*
Using the different reporters, we can run the tests, then
report the results.

Defaults to the JUnitTt reporter.

For inspecting per test result, VerboseTt.report may be useful
*)
structure Runner :> TestRunner = MkRunner(JUnitTt);

val tests : \nwlinkedidentc{Test.t}{NWsR2q7-3fh2dK-1} list = tl [
  \nwlinkedidentc{Test.suite}{NWsR2q7-3fh2dK-1} "" []
\LA{}Register unit tests~{\nwtagstyle{}\subpageref{nw@notdef}}\RA{}
];

fun main () = Runner.run tests;
\nwnotused{tests/main.sml}\nwidentuses{\\{{\nwixident{Test.suite}}{Test.suite}}\\{{\nwixident{Test.t}}{Test.t}}}\nwindexuse{\nwixident{Test.suite}}{Test.suite}{NWsR2q7-3MeQDa-1}\nwindexuse{\nwixident{Test.t}}{Test.t}{NWsR2q7-3MeQDa-1}\nwendcode{}\nwbegindocs{58}\nwdocspar

\begin{node}[SML/NJ adapter]
If we want to support SML/NJ, then we should organize test suites as
structures implementing the {\Tt{}TestSuite\nwendquote} signature, then add to
the {\Tt{}\LA{}Unit tests~{\nwtagstyle{}\subpageref{nw@notdef}}\RA{}\nwendquote} chunk the {\Tt{}MyTestSuite.suite\nwendquote}. We informally
stipulate that the {\Tt{}suite\nwendquote} value should be a test suite, but there's
no way to enforce this.
\end{node}

\nwenddocs{}\nwbegincode{59}\sublabel{NWsR2q7-1mlFrS-1}\nwmargintag{{\nwtagstyle{}\subpageref{NWsR2q7-1mlFrS-1}}}\moddef{tests/TestSuite.sig~{\nwtagstyle{}\subpageref{NWsR2q7-1mlFrS-1}}}\endmoddef\nwstartdeflinemarkup\nwenddeflinemarkup
signature TestSuite = sig
  val suite : \nwlinkedidentc{Test.t}{NWsR2q7-3fh2dK-1}
end;
\nwnotused{tests/TestSuite.sig}\nwidentuses{\\{{\nwixident{Test.t}}{Test.t}}}\nwindexuse{\nwixident{Test.t}}{Test.t}{NWsR2q7-1mlFrS-1}\nwendcode{}\nwbegindocs{60}\nwdocspar

\begin{node}[Build script for MLton]
Lastly, it will be convenient to have a build script for MLton (and
Poly/ML if the user employs \texttt{polymlb}).
\end{node}

\nwenddocs{}\nwbegincode{61}\sublabel{NWsR2q7-2v996w-1}\nwmargintag{{\nwtagstyle{}\subpageref{NWsR2q7-2v996w-1}}}\moddef{test.mlb~{\nwtagstyle{}\subpageref{NWsR2q7-2v996w-1}}}\endmoddef\nwstartdeflinemarkup\nwenddeflinemarkup
$(SML_LIB)/basis/basis.mlb
tests/Assert.sml
tests/Test.sig
tests/Test.sml
tests/Reporter.sig
tests/JUnitTt.sml
tests/TestRunner.sig
tests/MkRunner.sml
tests/TestSuite.sig
common.mlb
\LA{}Unit tests to run~{\nwtagstyle{}\subpageref{nw@notdef}}\RA{}
tests/main.sml
src/main.sml
\nwnotused{test.mlb}\nwendcode{}\nwbegindocs{62}\nwdocspar

\begin{node}[References]
Beck~\cite{beck1994simple} coded the progenitor SUnit framework for
treating specification documents as Smalltalk code.
Following Meszaros~\cite{meszaros2007xunit}, we describe the basic
design of a simplified testing framework for our needs.
\end{node}

\begin{xca}[Saving artifacts]
Extreme programming practices advocate test reporters saving an
``artifact'' of some kind when running the tests. Write a test
reporter producing some file recording the results.
\begin{enumerate}
\item Determine the markup for the file (XML? JSON? YAML?)
\item Write a specification for the artifact produced. 
\item Write a reporter producing this sort of artifact. Where will the artifact be placed? 
\item Write a tool which reads in the artifact, and prints out a summary to the screen. 
\end{enumerate}
\end{xca}

\begin{xca}[Combine test reporters]
Write a \SML/ functor \texttt{MkJointReporter(R1 : Reporter, R2 : Reporter) : Reporter}
which will perform all the reporting in {\Tt{}R1\nwendquote} and {\Tt{}R2\nwendquote}.
\end{xca}

\begin{xca}[Capture output]
Write code to capture \texttt{stdout} (and \texttt{stderr}), and store
anything written to these output streams as metadata in a test result.
\end{xca}

\begin{xca}[Exit status]
Modify the test runner code so it will exit with a failure status if
there were failed unit tests (or errors encountered).
\end{xca}

\begin{xca}[Testing the tester]
We mentioned that unit testing originated as a way to transform a
specification document into executable code.
\begin{enumerate}
\item Write specifications for the code in this section (i.e.,
  specifications for the unit testing framework itself)
\item Transform these specifications into unit tests, so we can check
  the unit testing code is doing what we expect.
\end{enumerate}
\end{xca}

\begin{xca}[BDD testing]
There are as many ways to test code as there are programmers (if not more).
A popular alternative to the minimal design we've implemented is
``behaviour driven development'' where you first write down the tests
in \define{Gherkin} to specify the behaviour expected in different
situations (``Given [setup], when [code is executed], expected
[outcome]'' is usually the template they use), then you run the tests
(to make sure they fail), then you implement the code to satisfy the
specified behaviour.
\begin{enumerate}
\item Research Gherkin and this approach to unit testing. What are its
  merits and drawbacks?
\item What would be needed to modify our code to implement a
  Gherkin-style testing framework?
\item Would this style of coding be well-suited to developing a proof
  assistant? Why or why not? (This could easily be an essay, a term
  paper, or a PhD thesis.)
\end{enumerate}
\end{xca}
\nwenddocs{}

\nwixlogsorted{c}{{Assert.sml}{NWsR2q7-4gQMdZ-1}{\nwixd{NWsR2q7-4gQMdZ-1}}}%
\nwixlogsorted{c}{{Constructors for tests}{NWsR2q7-3fh2dK-1}{\nwixu{NWsR2q7-3ucTCC-1}\nwixd{NWsR2q7-3fh2dK-1}}}%
\nwixlogsorted{c}{{Implementation of test result}{NWsR2q7-oQiQW-1}{\nwixu{NWsR2q7-3ucTCC-1}\nwixd{NWsR2q7-oQiQW-1}\nwixd{NWsR2q7-oQiQW-2}\nwixd{NWsR2q7-oQiQW-3}\nwixd{NWsR2q7-oQiQW-4}\nwixd{NWsR2q7-oQiQW-5}\nwixd{NWsR2q7-oQiQW-6}\nwixd{NWsR2q7-oQiQW-7}\nwixd{NWsR2q7-oQiQW-8}\nwixd{NWsR2q7-oQiQW-9}}}%
\nwixlogsorted{c}{{Iteratively report JUnit-style report for a single outcome}{NWsR2q7-4wVoS-1}{\nwixu{NWsR2q7-bqErk-1}\nwixd{NWsR2q7-4wVoS-1}}}%
\nwixlogsorted{c}{{JUnitTt.sml}{NWsR2q7-bqErk-1}{\nwixd{NWsR2q7-bqErk-1}}}%
\nwixlogsorted{c}{{Register unit tests}{nw@notdef}{\nwixu{NWsR2q7-3MeQDa-1}}}%
\nwixlogsorted{c}{{Reporter.sig}{NWsR2q7-jZez2-1}{\nwixd{NWsR2q7-jZez2-1}}}%
\nwixlogsorted{c}{{Run a test case or suite}{NWsR2q7-1bpV8D-1}{\nwixu{NWsR2q7-3ucTCC-1}\nwixd{NWsR2q7-1bpV8D-1}}}%
\nwixlogsorted{c}{{Specification for test result}{NWsR2q7-1ETFk0-1}{\nwixu{NWsR2q7-1Fzsdc-1}\nwixd{NWsR2q7-1ETFk0-1}\nwixd{NWsR2q7-1ETFk0-2}\nwixd{NWsR2q7-1ETFk0-3}\nwixd{NWsR2q7-1ETFk0-4}\nwixd{NWsR2q7-1ETFk0-5}\nwixd{NWsR2q7-1ETFk0-6}\nwixd{NWsR2q7-1ETFk0-7}\nwixd{NWsR2q7-1ETFk0-8}\nwixd{NWsR2q7-1ETFk0-9}}}%
\nwixlogsorted{c}{{test.mlb}{NWsR2q7-2v996w-1}{\nwixd{NWsR2q7-2v996w-1}}}%
\nwixlogsorted{c}{{Test.sig}{NWsR2q7-1Fzsdc-1}{\nwixd{NWsR2q7-1Fzsdc-1}}}%
\nwixlogsorted{c}{{Test.sml}{NWsR2q7-3ucTCC-1}{\nwixd{NWsR2q7-3ucTCC-1}}}%
\nwixlogsorted{c}{{TestRunner.sig}{NWsR2q7-4Peuwc-1}{\nwixd{NWsR2q7-4Peuwc-1}}}%
\nwixlogsorted{c}{{tests/main.sml}{NWsR2q7-3MeQDa-1}{\nwixd{NWsR2q7-3MeQDa-1}}}%
\nwixlogsorted{c}{{tests/MkRunner.sml}{NWsR2q7-2dKvEw-1}{\nwixd{NWsR2q7-2dKvEw-1}}}%
\nwixlogsorted{c}{{tests/TestSuite.sig}{NWsR2q7-1mlFrS-1}{\nwixd{NWsR2q7-1mlFrS-1}}}%
\nwixlogsorted{c}{{Unit tests to run}{nw@notdef}{\nwixu{NWsR2q7-2v996w-1}}}%
\nwixlogsorted{i}{{\nwixident{Assert.eq}}{Assert.eq}}%
\nwixlogsorted{i}{{\nwixident{Assert.Fail}}{Assert.Fail}}%
\nwixlogsorted{i}{{\nwixident{Assert.that}}{Assert.that}}%
\nwixlogsorted{i}{{\nwixident{Test.mk}}{Test.mk}}%
\nwixlogsorted{i}{{\nwixident{Test.Result.count{\_}errors}}{Test.Result.count:unerrors}}%
\nwixlogsorted{i}{{\nwixident{Test.Result.count{\_}failures}}{Test.Result.count:unfailures}}%
\nwixlogsorted{i}{{\nwixident{Test.Result.count{\_}successes}}{Test.Result.count:unsuccesses}}%
\nwixlogsorted{i}{{\nwixident{Test.Result.count{\_}total}}{Test.Result.count:untotal}}%
\nwixlogsorted{i}{{\nwixident{Test.Result.exn{\_}message}}{Test.Result.exn:unmessage}}%
\nwixlogsorted{i}{{\nwixident{Test.Result.for{\_}case}}{Test.Result.for:uncase}}%
\nwixlogsorted{i}{{\nwixident{Test.Result.for{\_}suite}}{Test.Result.for:unsuite}}%
\nwixlogsorted{i}{{\nwixident{Test.Result.is{\_}case}}{Test.Result.is:uncase}}%
\nwixlogsorted{i}{{\nwixident{Test.Result.is{\_}error}}{Test.Result.is:unerror}}%
\nwixlogsorted{i}{{\nwixident{Test.Result.is{\_}failure}}{Test.Result.is:unfailure}}%
\nwixlogsorted{i}{{\nwixident{Test.Result.is{\_}success}}{Test.Result.is:unsuccess}}%
\nwixlogsorted{i}{{\nwixident{Test.Result.is{\_}suite}}{Test.Result.is:unsuite}}%
\nwixlogsorted{i}{{\nwixident{Test.Result.msg}}{Test.Result.msg}}%
\nwixlogsorted{i}{{\nwixident{Test.Result.name}}{Test.Result.name}}%
\nwixlogsorted{i}{{\nwixident{Test.Result.realtime}}{Test.Result.realtime}}%
\nwixlogsorted{i}{{\nwixident{Test.Result.results}}{Test.Result.results}}%
\nwixlogsorted{i}{{\nwixident{Test.Result.runtime}}{Test.Result.runtime}}%
\nwixlogsorted{i}{{\nwixident{Test.Result.t}}{Test.Result.t}}%
\nwixlogsorted{i}{{\nwixident{Test.run}}{Test.run}}%
\nwixlogsorted{i}{{\nwixident{Test.suite}}{Test.suite}}%
\nwixlogsorted{i}{{\nwixident{Test.t}}{Test.t}}%

